\documentclass[12pt,oneside,a4paper,brazil]{abntex2}
\usepackage[T1]{fontenc}
\usepackage[utf8]{inputenc}
\usepackage[scaled=0.95]{helvet}
\renewcommand{\familydefault}{\sfdefault}
\usepackage{microtype}
\usepackage{indentfirst}
\usepackage{graphicx}
\usepackage{enumitem}
\usepackage{xurl}
\usepackage{hyperref}
\usepackage[numbers]{natbib}
\usepackage{setspace}
\usepackage[top=3cm,left=3cm,right=2cm,bottom=2cm]{geometry}
\hypersetup{colorlinks=true,linkcolor=black,citecolor=black,urlcolor=blue,
 pdftitle={Planeta Consciente: desenvolvimento de um site estático interativo para educação ambiental},
 pdfauthor={Ana Santiago},pdfsubject={Artigo de Projeto Integrador — Curso Técnico em Análise e Desenvolvimento de Sistemas}}
\setlength{\parindent}{1.25cm}
\setlength{\parskip}{0.15cm}
\OnehalfSpacing
\urlstyle{same}
\newcommand{\projectversion}{0.3.0-projeto.1}
\newcommand{\tituloartigo}{Planeta Consciente: desenvolvimento de um site estático interativo para educação ambiental}
\begin{document}
\begin{capa}
\begin{center}
\textbf{CENTRO DE EDUCAÇÃO DE TEMPO INTEGRAL (CETI) LUCAS MEIRELES}\\
\textbf{CURSO TÉCNICO EM ANÁLISE E DESENVOLVIMENTO DE SISTEMAS}\\[0.4cm]
Código INEP 22161201 — localização rural\\[1.8cm]
\textbf{PROJETO INTEGRADOR}\\[2.4cm]
\textbf{ANA SANTIAGO}\\[2.4cm]
\textbf{\MakeUppercase{\tituloartigo}}\\[2.4cm]
BR-316, km 24, s/n, Chapadinha Sul\\
Teresina — Piauí\\2026
\end{center}
\vfill
\begin{center}\textbf{Orientador: Aloisio de Araújo Costa Magalhães}\end{center}
\end{capa}

\begin{folhaderosto}
\begin{center}\textbf{ANA SANTIAGO}\end{center}
\vspace{2.5cm}
\begin{center}\textbf{\MakeUppercase{\tituloartigo}}\end{center}
\vspace{1.5cm}
\hfill\begin{minipage}{0.52\textwidth}
Artigo apresentado ao Curso Técnico em Análise e Desenvolvimento de Sistemas do Centro de Educação de Tempo Integral Lucas Meireles como atividade do Projeto Integrador.\\[0.5cm]
Orientador: Aloisio de Araújo Costa Magalhães.
\end{minipage}
\vfill
\begin{center}Teresina — Piauí\\2026\end{center}
\end{folhaderosto}

\begin{resumo}
Este artigo relata o desenvolvimento do Planeta Consciente, um site estático interativo de educação ambiental produzido como Projeto Integrador do Curso Técnico em Análise e Desenvolvimento de Sistemas do CETI Lucas Meireles. O objetivo foi organizar conteúdos temáticos, gráficos, ações cotidianas, curiosidades, notícias e um mapa interativo de biomas em uma aplicação web responsiva e publicável sem servidor de aplicação. O trabalho adotou um processo descritivo de engenharia: levantamento do escopo, organização dos dados, implementação modular em Next.js e React, exportação estática, integração contínua e publicação no GitHub Pages. O artefato reúne navegação por seções, gráficos renderizados com Recharts, seleção de estados em mapa SVG e controles interativos de interface. A verificação realizada foi técnica: o build e a implantação automatizada foram concluídos pelo GitHub Actions. Não foram conduzidos testes com participantes nem avaliação formal de aprendizagem, usabilidade ou acessibilidade. A inspeção do código identificou limitações de rastreabilidade das fontes ambientais, simplificação da associação estado-bioma e necessidade de revisar governança da ferramenta analítica carregada em produção. Conclui-se que o projeto entregou um protótipo funcional publicado, cuja evolução depende de validação factual, acessibilidade e avaliação com usuários.

\vspace{\onelineskip}
\noindent\textbf{Palavras-chave}: educação ambiental; desenvolvimento web; Next.js; visualização de dados; mapas interativos.
\end{resumo}

\noindent\textbf{Versão do artigo}: \projectversion.\\
\pdfbookmark[0]{\contentsname}{toc}
\tableofcontents*
\textual
\chapter{Introdução}\label{page:introducao:start}
A educação ambiental em meios digitais pode combinar explicações, dados e interação para apresentar temas como água, florestas, reciclagem, mudanças climáticas e biomas. O desafio de desenvolvimento é organizar esses recursos em uma interface navegável sem confundir visualização atraente com informação verificada. Este trabalho apresenta o Planeta Consciente como artefato de software construído no Projeto Integrador do Curso Técnico em Análise e Desenvolvimento de Sistemas.

O problema de projeto foi estruturar, em um único site, conteúdos ambientais e componentes interativos que pudessem ser publicados em hospedagem estática. A questão orientadora é: como implementar e publicar uma aplicação web responsiva que organize conteúdos ambientais, gráficos e um mapa interativo, mantendo uma arquitetura simples e passível de evolução? O escopo deste artigo é o desenvolvimento do artefato e a verificação técnica de compilação e publicação; não se afirma que o site aumente conhecimento, mude comportamento ou tenha usabilidade validada por participantes.

A escolha por exportação estática viabiliza servir HTML, CSS, JavaScript e arquivos públicos sem manter um servidor Next.js em execução, respeitadas as limitações desse modo de implantação \cite{nextstatic}. A aplicação foi publicada no GitHub Pages por fluxo de integração contínua, que constrói o site e envia os arquivos estáticos ao ambiente de publicação \cite{ghpages}.
\label{page:introducao:end}

\chapter{Objetivos}\label{page:objetivos:start}
\section{Objetivo geral}
Desenvolver e publicar o Planeta Consciente como site estático interativo de educação ambiental, organizando conteúdos, gráficos e mapa em uma aplicação web modular e verificável.

\section{Objetivos específicos}
\begin{enumerate}[leftmargin=1.25cm]
\item Implementar uma interface responsiva composta por seções temáticas, gráficos, ações sustentáveis, curiosidades, notícias e mapa de biomas, com dados estruturados e componentes reutilizáveis.
\item Automatizar a geração e a publicação do site e dos artefatos acadêmicos, registrando as verificações técnicas executadas e as limitações que requerem validação futura.
\end{enumerate}
\label{page:objetivos:end}

\chapter{Metodologia}\label{page:metodologia:start}
Este trabalho é um relato descritivo de desenvolvimento de software, delimitado ao repositório versionado do Planeta Consciente. A unidade de análise é a aplicação e seus componentes, seus dados locais e seus fluxos automatizados de construção e publicação. O procedimento foi organizado em quatro etapas: delimitação do escopo; decomposição da página em componentes; implementação e integração dos dados; e verificação por build de produção e implantação estática. A inspeção é baseada no código e nos workflows do repositório; não corresponde a experimento nem a avaliação com usuários.

A página principal compõe cabeçalho, apresentação, conscientização, temas, atitudes sustentáveis, curiosidades, mapa, notícias e rodapé. A camada de conteúdo em TypeScript separa dados dos componentes de apresentação. Next.js com App Router e React fornece a composição da interface; Recharts é usado para gráficos, enquanto o mapa consome geometrias SVG locais e uma tabela de associação entre siglas estaduais e biomas. A aplicação também utiliza Tailwind CSS para estilos e Lucide para ícones. A descrição das bibliotecas segue sua documentação primária \cite{reactstate,recharts}.

Na verificação técnica, o fluxo de integração instala dependências com arquivo de lock, executa o build estático, valida a estrutura do manuscrito e compila PDFs com LaTeX. O fluxo de publicação transfere a pasta exportada para GitHub Pages. A Release versionada empacota o site e os documentos compilados. Esses procedimentos verificam geração e entrega do artefato; não verificam correção científica de cada dado ambiental, compatibilidade com tecnologias assistivas ou experiência de pessoas usuárias.
\label{page:metodologia:end}

\chapter{Referencial teórico}\label{page:referencial:start}
\section{Engenharia de aplicações web estáticas}
A decomposição de interface em componentes permite manter partes com responsabilidades distintas e integrar conteúdo, estado e apresentação. React descreve estado como mecanismo para refletir entradas e mudanças de interface, o que fundamenta controles locais como seleção de tema, estado no mapa e marcação de atitudes \cite{reactstate}. No projeto, o conteúdo temático e os biomas são mantidos em módulos de dados TypeScript, e a página coordena componentes especializados.

A exportação estática do Next.js gera arquivos que podem ser servidos por hospedagem estática; recursos que dependem de servidor em tempo de requisição ficam fora desse modelo \cite{nextstatic}. Essa propriedade é compatível com GitHub Pages, cujo fluxo oficial de Actions constrói o site, publica o artefato e registra a URL do ambiente \cite{ghpages}. A decisão reduz a necessidade de infraestrutura de execução, mas implica que atualizações do conteúdo dependem de nova compilação e publicação.

\section{Visualização e rastreabilidade de informação}
Gráficos e mapas são interfaces de interpretação, não apenas elementos decorativos. Uma visualização exige contexto para que seu público compreenda unidade, período, escala, fonte e limites do dado. Recharts oferece componentes composáveis de gráficos para React; a biblioteca implementa a camada de desenho, mas não valida a veracidade nem a adequação semântica dos valores fornecidos \cite{recharts}. Por isso, cada afirmação ambiental e cada série precisam de fonte específica, recorte e método documentados.

No mapa, a associação de um único bioma predominante a cada estado simplifica uma realidade em que biomas e transições não coincidem necessariamente com fronteiras estaduais. O recurso serve como exploração introdutória e não deve ser apresentado como cartografia oficial ou substituto dos limites temáticos do IBGE \cite{ibge}. A qualidade científica do conteúdo depende de uma auditoria de afirmações e fontes que ainda não foi concluída.

\section{Integridade e apresentação acadêmica}
A redação deste artigo separa funcionalidades efetivamente implementadas de avaliações ainda não realizadas. Fontes técnicas são citadas junto às afirmações que sustentam, e as referências incluem URLs rastreáveis. A versão final deverá seguir o manual institucional e as normas ABNT aplicáveis à apresentação, às citações e às referências; o uso de uma classe LaTeX não comprova conformidade automática. As diretrizes do CNPq reforçam a responsabilidade de atribuição e integridade das fontes \cite{cnpq}.
\label{page:referencial:end}

\chapter{Resultados e discussão}
O artefato implementado reúne uma página única com navegação por âncoras e sete áreas de conteúdo. A seção de temas seleciona tópicos e apresenta séries em gráficos de barras, linhas ou setores; a área de atitudes atualiza progresso enquanto a página está aberta; curiosidades usam controles expansíveis; e o mapa permite selecionar estado e filtrar biomas. As notícias são blocos editoriais estáticos, sem integração com serviço externo de notícias. Essas funções foram confirmadas na composição da página e nos componentes e dados do código-fonte.

A exportação estática foi construída e publicada automaticamente no endereço do projeto. A verificação remota do GitHub Actions confirmou a conclusão bem-sucedida do build de produção e da publicação no Pages. A Release \texttt{v0.2.0-protocolo.7} também disponibilizou ZIP do site e PDFs acadêmicos; o novo artigo será incorporado à próxima versão. Esses resultados demonstram entrega e reprodutibilidade do processo de build, não qualidade de conteúdo ou impacto educacional.

A inspeção do artefato identificou quatro limites relevantes. Primeiro, alguns dados e notas de gráficos mencionam fontes em termos gerais, sem ligação individual e metadados completos por afirmação. Segundo, o mapa usa uma categoria predominante por estado, portanto comunica uma simplificação espacial. Terceiro, as atitudes sustentáveis são armazenadas apenas no estado da interface e não persistem após recarregamento. Quarto, o layout inclui a ferramenta Vercel Analytics em produção; antes de uso educacional ampliado, é necessário documentar sua configuração, os dados tratados e as informações de privacidade conforme a LGPD \cite{lgpd}. Nenhum desses pontos foi resolvido por testes com usuários.

Também permanece uma divergência cadastral a conferir: os dados fornecidos para a capa informam BR-316, km 24 e CEP 64022-990, enquanto publicação do Diário Oficial do Piauí associa o código INEP 22161201 ao CETI Lucas Meireles, zona rural de Teresina, com endereço em km 22 e CEP 64390-000 \cite{doepi2025}. O nome institucional foi associado ao código; endereço e CEP devem ser confirmados com a escola antes da entrega definitiva.
\label{page:resultados:end}

\chapter{Considerações finais}\label{page:consideracoes:start}
O Projeto Integrador resultou no desenvolvimento e na publicação do Planeta Consciente, uma aplicação estática que reúne conteúdos temáticos, gráficos, curiosidades, ações sustentáveis, notícias e um mapa interativo. A arquitetura baseada em Next.js, React, módulos de dados e componentes especializados permitiu gerar os arquivos estáticos e automatizar sua publicação no GitHub Pages. Os objetivos de implementação e entrega técnica foram atendidos conforme as verificações de build e implantação executadas.

O resultado não autoriza concluir que o site melhora aprendizagem ou comportamento ambiental: não houve estudo com participantes, teste de usabilidade, avaliação formal de acessibilidade ou medição de desempenho educacional. Para evoluir o projeto, recomenda-se vincular cada dado ambiental a fonte verificável, rever simplificações do mapa, avaliar acessibilidade com critérios próprios, definir transparência e governança de Analytics e realizar testes de experiência com autorização e salvaguardas adequadas. O endereço escolar divergente também deve ser confirmado antes do protocolo final. Assim, o produto é uma entrega técnica funcional com agenda explícita de validação factual e de experiência de uso.
\label{page:consideracoes:end}

\postextual
\begin{thebibliography}{99}
\label{page:referencias:start}
\bibitem{nextstatic} VERCEL. \textit{Static exports}. Next.js Documentation. Disponível em: \url{https://nextjs.org/docs/app/guides/static-exports}. Acesso em: 6 out. 2026.
\bibitem{reactstate} META OPEN SOURCE. \textit{Managing state}. React Documentation. Disponível em: \url{https://react.dev/learn/managing-state}. Acesso em: 6 out. 2026.
\bibitem{recharts} RECHARTS GROUP. \textit{Recharts: composable charting library built on React components}. Disponível em: \url{https://recharts.org/}. Acesso em: 6 out. 2026.
\bibitem{ghpages} GITHUB. \textit{Deploying your website automatically}. GitHub Docs. Disponível em: \url{https://docs.github.com/en/get-started/start-your-journey/deploying-your-website-automatically}. Acesso em: 6 out. 2026.
\bibitem{ibge} INSTITUTO BRASILEIRO DE GEOGRAFIA E ESTATÍSTICA. \textit{Biomas e Sistema Costeiro-Marinho do Brasil}. Disponível em: \url{https://www.ibge.gov.br/geociencias/informacoes-ambientais/vegetacao/15842-biomas.html?lang=pt-BR}. Acesso em: 6 out. 2026.
\bibitem{doepi2025} PIAUÍ. Diário Oficial do Estado. \textit{Portaria sobre alteração de vinculação administrativa do CETI Lucas Meireles}, publicada em 23 jul. 2025. Disponível em: \url{https://www.diario.pi.gov.br/doe/files/diarios/anexo/717cef99-df56-4a1c-81e1-e21f586dbf0e/DOEPI_139_2025.pdf}. Acesso em: 6 out. 2026.
\bibitem{cnpq} BRASIL. Conselho Nacional de Desenvolvimento Científico e Tecnológico. \textit{Diretrizes de integridade na pesquisa apoiada pelo CNPq}. Disponível em: \url{https://www.gov.br/cnpq/pt-br/composicao/comissao-de-integridade/diretrizes}. Acesso em: 6 out. 2026.
\bibitem{lgpd} BRASIL. \textit{Lei nº 13.709, de 14 de agosto de 2018: Lei Geral de Proteção de Dados Pessoais}. Texto compilado. Disponível em: \url{https://www.planalto.gov.br/ccivil_03/_ato2015-2018/2018/lei/l13709compilado.htm}. Acesso em: 6 out. 2026.
\end{thebibliography}
\label{page:referencias:end}
\end{document}
